\documentclass[a4paper,12pt]{article} % Тип документа (article можно поменять на book при желании)
\usepackage[utf8]{inputenc} % Кодировка исходного текста
\usepackage[T2A]{fontenc} % Кодировка результирующего текста
\usepackage[english,russian]{babel} % Подключение русского языка 
\usepackage[left=25mm, top=20mm, right=25mm, bottom=30mm, nohead, nofoot]{geometry} 
\usepackage{amsmath,amsfonts,amssymb,amsthm} % Математика
\usepackage{fancybox,fancyhdr}
\usepackage{graphicx} % Графика
\usepackage[dvipsnames]{xcolor} % Для разноцветного текста
\usepackage{eso-pic,graphicx} % Для создания фоновой картинки
\usepackage{anyfontsize} % Для разблокировки всех размеров текста
\usepackage[colorlinks=true, linkcolor=WildStrawberry]{hyperref} % Для гиперссылок
\usepackage{float}
\usepackage[many]{tcolorbox}
\usepackage[document]{ragged2e}
\usepackage{wrapfig} % Для обтекания
\usepackage{enumitem} % для улучшенных списков
\usepackage{multicol} %для разделения на колонны.
\usepackage{circuitikz}
\usepackage{longtable} %для длинных таблиц

% Обьявление команд:
% Объявляем окружение "definition"
\newtheorem{definition}{Определение} % [section] - опционально для нумерации внутри раздела
\newtheorem*{notice}{Замечание}
\newtheorem{theorem}{Теорема}
\newtheorem{statement}{Утверждение}
\newcommand{\bbox}[3][]


\begin{document}

% command to provide stretchy vertical space in proportion
\newcommand\nbvspace[1][3]{\vspace*{\stretch{#1}}}
% allow some slack to avoid under/overfull boxes
\newcommand\nbstretchyspace{\spaceskip0.5em plus 0.25em minus 0.25em}
% To improve spacing on titlepages
\newcommand{\nbtitlestretch}{\spaceskip0.6em}
\pagestyle{empty}
\begin{center}
\bfseries
\nbvspace[1]
\Huge
{\fontsize{44}{60}\selectfont
\textcolor{black}{Теория вероятности. Статистика и Анализ распределений}}

\nbvspace[1]
\huge \textcolor{Maroon}{Институт Интеллектуальных Кибернетических Систем}\\[0.5em]

\nbvspace[2]


\nbvspace[3]
\normalsize

\textcolor{black}{Москва, НИЯУ "МИФИ"\space}\\
\large
\textcolor{black}{Осень 2026}
\nbvspace[1]
\end{center}


%новая страница


\newpage

\pagestyle{fancy}
\fancyhf{}
\fancyhead[R]{\href{https://youtube.com/playlist?list=PLUWSgPsx43FOzb0GhQMH998Fpw99t3OcG&si=YlHxivoaM8SGhgkg}{\textcolor{Maroon}{ИИКС ПМИ}}}
\fancyfoot[R]{\thepage} 
\fancyhead[L]{Статистика и Анализ распределений}
\setcounter{page}{2}
\headsep=10mm

\tableofcontents

%новая страница

\newpage


\section*{Введение}
\addcontentsline{toc}{section}{Введение}

Теория вероятностей представляет собой фундаментальный раздел математики, изучающий закономерности не просто в отдельных случайных явлениях, а именно в \textbf{массовых случайных явлениях}. Только при наличии достаточно большой массы наблюдений, проводимых в одних и тех же условиях, становится возможным установление устойчивых закономерностей в случайных процессах. 

Рассмотрим классический пример: бросание монеты. Выпадение «герба» или «решки» в результате единичного броска является событием чисто случайным. Теория вероятностей не ставит своей целью объяснить, почему в данном конкретном случае выпал «герб», а не «решка». Ее фундаментальная задача заключается в ином: определить, насколько вероятно выпадение «герба» при условии, что монета будет подброшена достаточно большое количество раз. Именно в пределе большого числа испытаний хаотичность отдельных событий уступает место строгим статистическим закономерностям.

К созданию и развитию теории вероятностей и математической статистики приложили свои усилия многие выдающиеся ученые. Среди пионеров и классиков этой науки следует назвать Г.~Галилея, Б.~Паскаля, П.~Ферма, Х.~Гюйгенса, Я.~Бернулли, А.~Муавра, П.~Лапласа, С.~Пуассона, а также ученых, внесших вклад в ее современное развитие: У.~Госсета (Стьюдента), К.~Пирсона, Н.~Винера, У.~Феллера, Дж.~Дуба, Р.~Фишера и многих других.

Особый вклад в становление и строгое аксиоматическое обоснование теории вероятностей внесли наши соотечественники. Славная плеяда российских и советских математиков включает В.~Я.~Буняковского, П.~Л.~Чебышева, А.~М.~Ляпунова, А.~А.~Маркова, С.~Н.~Бернштейна, А.~Я.~Хинчина, А.~Н.~Колмогорова, В.~И.~Романовского, Н.~В.~Смирнова и многих других исследователей, чьи работы заложили незыблемый фундамент современной вероятностно-статистической науки.

\vspace{1em}
\noindent\textit{Данная методичка призвана помочь вам систематизировать знания по этим темам, освоить ключевые методы анализа распределений и научиться применять их для решения практических задач.}

\newpage

\section*{Обозначения}
\addcontentsline{toc}{section}{Обозначения}

\begin{longtable}{|p{3cm}|p{4cm}|p{3cm}|p{4cm}|}
\hline
\textbf{Обозначение} & \textbf{Определение} & \textbf{Обозначение} & \textbf{Определение} \\
\hline
\endfirsthead

\hline
\textbf{Обозначение} & \textbf{Определение} & \textbf{Обозначение} & \textbf{Определение} \\
\hline
\endhead

\hline
\endfoot

\hline
\endlastfoot

$\mathfrak{A}$ & алгебра множеств & $o(\cdot)$ & о-малое \\
\hline
$\mathrm{B}(n, p)$ & биномиальное распределение & $|A|$ & определитель матрицы \\
\hline
$\mathfrak{B}(\cdot)$ & борелевская $\sigma$-алгебра над $X$ & $A \cap B$ & пересечение множеств/событий \\
\hline
$\mathcal{P}(\cdot)$ & булеан & $AB$ & пересечение событий \\
\hline
$\mathbb{P}$ & вероятность, вероятностная мера & $f_{\xi}$ & плотность распределения \\
\hline
$\mu^*$ & внешняя мера & $\mathcal{F}\{\cdot\}$ & преобразование Фурье \\
\hline
$\mu_*$ & внутренняя мера & $\Omega$ & пространство элементарных событий \\
\hline
$\Gamma(x)$ & гамма-функция & $\varnothing$ & пустое множество \\
\hline
$\Gamma(k, \theta)$ & гамма-распределение & $A \setminus B$ & разность множеств/событий \\
\hline
$\mathrm{Geom}(p)$ & геометрическое распределение & $\mathcal{U}(a, b)$ & равномерное распределение \\
\hline
$\mathrm{HG}(D, N, n)$ & гипергеометрическое распределение & $\mathrm{Ber}(p)$ & распределение Бернулли \\
\hline
$A^n$ & декартова степень & $\mathrm{Cauchy}(x_0, \gamma)$ & распределение Коши \\
\hline
$A \times B$ & декартово произведение & $\mathrm{P}(\lambda)$ & распределение Пуассона \\
\hline
$\mathrm{diag}(\cdot, ..., \cdot)$ & диагональная матрица & $\chi^2(n)$ & распределение хи-квадрат \\
\hline
$A \sqcup B$ & дизъюнктное объединение множеств/событий & $F_{\xi|\eta}$ & регулярная функция распределения \\
\hline
$\mathbb{D}\xi$ & дисперсия & $\mathfrak{S}$ & $\sigma$-алгебра \\
\hline
$\overline{A}$ & дополнение события & $A \triangle B$ & симметрическая разность множеств \\
\hline
$\mathbf{1}_A(x)$ & индикатор множества/события & $(\vec{a}, \vec{b})$ & скалярное произведение \\
\hline
$\int_a^b f(x)\,\mathrm{d}x$ & интеграл Римана & $\vec{\xi}$ & случайный вектор \\
\hline
$\int_X f(x) \mu(\mathrm{d}x)$ & интеграл Лебега & $\xi$ & случайная величина \\
\hline
$\mathbf{1}_A(x)$ & индикатор множества/события & $a \stackrel{d}{=} b$ & сравнимость по модулю \\
\hline
$\mathrm{cov}(\cdot)$ & ковариационная матрица & $\sigma$ & стандартное отклонение \\
\hline
$\mathrm{cov}(\cdot, \cdot)$ & ковариация & $A \subset B$ & строгое подмножество \\
\hline
$A_s$ & коэффициент асимметрии & $\xrightarrow{L^p}$ & сходимость в смысле $L^p$ \\
\hline
$\rho(\cdot, \cdot)$ & коэффициент корреляции & $\xrightarrow{\mathbb{P}}$ & сходимость по вероятности \\
\hline
$\gamma_2$ & коэффициент эксцесса & $\xrightarrow{d}$ & сходимость по распределению \\
\hline
$\nu_k$ & $k$-й начальный момент & $\xrightarrow{\text{п.н.}}$ & сходимость почти наверное \\
\hline
$\mu_k$ & $k$-й центральный момент & $\sup$ & точная верхняя грань \\
\hline
$L(\mu, s)$ & логистическое распределение & $\inf$ & точная нижняя грань \\
\hline
$\mathbb{E}\xi$ & математическое ожидание & $F_{\xi}$ & функция распределений \\
\hline
$\mu$ & мера & $A^T$ & транспонированная матрица \\
\hline
$\mathfrak{S}(X)$ & минимальная $\sigma$-алгебра & $\mathbb{P}(A | \cdot)$ & условная вероятность \\
\hline
$\mathcal{N}(\vec{\mu}, \Sigma)$ & многомерное нормальное распределение & $\mathbb{E}[\xi | \cdot]$ & условное мат. ожидание \\
\hline
$\mathbb{R}$ & множество действительных чисел & $f_{\xi|\eta}$ & условная плотность распределения \\
\hline
$\mathbb{C}$ & множество комплексных чисел & $n!$ & факториал числа \\
\hline
$\mathbb{N}$ & множество натуральных чисел & $\mathrm{erf}(\cdot)$ & функция ошибок \\
\hline
$A \subseteq B$ & нестрогое подмножество & $\varphi_{\xi}$ & характеристическая функция \\
\hline
$\mathcal{N}(\mu, \sigma^2)$ & нормальное распределение & $P_n$ & число перестановок \\
\hline
$O(\cdot)$ & о-большое & $A_n^k$ & число размещений из $n$ по $k$ \\
\hline
$A^{-1}$ & обратная матрица & $C_n^k$ & число сочетаний из $n$ по $k$ \\
\hline
$f^{-1}$ & обратная функция & $\mathrm{Exp}(\lambda)$ & экспоненциальное распределение \\
\hline
$A \cup B$ & объединение множеств/событий & & \\
\hline

\end{longtable}

\newpage
\section{Алгебра событий}
\subsection{Основные определения}
\subsubsection{События. Исходы. Пространство элементарных исходов}

Для начала следует познакомиться с тремя основополагающими понятиями в Алгебре событий:

\begin{definition}[Элементарный исход]
\textbf{Элементарный исход} (элементарное событие, элементарный исход эксперимента) -- это отдельный, далее неделимый результат случайного эксперимента, обозначаемый как $\omega \in \Omega$. Каждый элементарный исход представляет собой один конкретный вариант того, чем может закончиться эксперимент.
\end{definition}

\begin{definition}[Событие]
\textbf{Событие} -- множество, состоящее из элементарных исходов данного эксперимента, т.е. некоторое подмножество пространства элементарных исходов: $A \subseteq \Omega$. Событие происходит (наступает), если в результате эксперимента реализуется хотя бы один элементарный исход $\omega$, принадлежащий этому множеству.
\end{definition}

\begin{definition}[Пространство элементарных исходов]
Пространство элементарных исходов/событий -- множество $\Omega$ событий (всех возможных исходов), являющихся исходами данного случайного эксперимента, из которых в эксперименте происходит ровно один.
\end{definition}

\textbf{Пример:}
    \center\textit{Эксперимент:} бросание игральной кости.
    \begin{itemize}
    \item Пространство элементарных исходов: $\Omega = \{1, 2, 3, 4, 5, 6\}$.
    
    \item Элементарный исход: $\omega = 4$ (выпала четвёрка).
    
    \item Событие $A$: ``выпало число больше 4'' $\Rightarrow A = \{5, 6\}$.
    
    \item Событие $B$: ``выпало нечётное число'' $\Rightarrow B = \{1, 3, 5\}$.
    \end{itemize}

\begin{notice}Таким образом:

\begin{enumerate}
    \item Элементарный исход -- это \textit{элемент} множества $\Omega$ (обозначается $\omega \in \Omega$).
    \item Событие -- это \textit{подмножество} $\Omega$ (обозначается $A \subseteq \Omega$).
    \item Событие может состоять из одного исхода (такое событие называется \textit{элементарным}) или из нескольких исходов.
    \item Всё пространство $\Omega$ тоже является событием (достоверное событие), равно как и пустое множество $\varnothing$ -- тоже событие (невозможное событие).
\end{enumerate}
\end{notice}

Элементарные исходы, которые являеются элементами данного события,
называют \textbf{благоприятными}.

\subsubsection{Операции над событиями}

Поскольку события мы определили как множества, то событиям свойственны множественные операции

\begin{definition}[Пересечение (произведение) событий]
Событие $C$, происходящее тогда и только тогда, когда наступает и событие $A$ и событие $B$, называется \textbf{пересечением (произведением) событий $A$ и $B$}. Обозначение: $A \cap B$ или $AB$.
\end{definition}

\begin{definition}[Несовместные события]
События $A$ и $B$ называют \textbf{несовместными}, если они не могут произойти одновременно, т.е. $A \cap B = \varnothing$.
\end{definition}

В противном случае события называют совместными или пересекающимися.

\begin{definition}[Объединение (сумма) событий]
Событие $C$, происходящее тогда и только тогда, когда происходит хотя бы одно из событий $A$ или $B$, называется \textbf{объединением (суммой) событий $A$ и $B$}. Обозначение: $A \cup B$ или $A + B$.
\end{definition}
\subsection{Классическая схема теории вероятностей}
Представьте, что вы играете в настольную игру и кидаете обычный игральный кубик. Каков шанс, что выпадет шестерка? Интуитивно понятно: один шанс из шести.
Это и есть "классическая схема теории вероятностей" (или классическое определение вероятности). Это самый простой, базовый и «честный» способ посчитать шанс наступления какого-то события. Если строго:
\begin{definition}[Классическая схема теории вероятностей]
Классическая схема теории вероятностей -- предположения, касаемые исследуемых экспериментов:
\begin{itemize}
    \item пространство элементарных исходов $\Omega$ -- конечное множество;
    \item элементарные исходы равновозможны.
\end{itemize}
\end{definition}

Последний пункт означает, что вероятность одинакова для каждого элементарного исхода. Тогда для событий:

\begin{definition}[Вероятность случайного события]
Вероятность случайного события $A$ в классической схеме -- отношение количества благоприятных исходов $N_A$ для события $A$ к общему числу исходов $N$:
\[
\mathbb{P}(A) = \frac{N_A}{N}.
\]
\end{definition}

Выделим довольно очевидные свойства:
\begin{enumerate}
    \item $\mathbb{P}(\varnothing) = 0$ и $\mathbb{P}(\Omega) = 1$;
    \item $\mathbb{P}(A \sqcup B) = \mathbb{P}(A) + \mathbb{P}(B)$;
    \item $\mathbb{P}(A \cup B) = \mathbb{P}(A) + \mathbb{P}(B) - \mathbb{P}(A \cap B)$;
    \item $\mathbb{P}(\overline{A}) = 1 - \mathbb{P}(\overline{B})$.
\end{enumerate}

Все задачи в классической схеме теории вероятностей сводятся так или иначе к комбинаторике, формулы которой будем считать известными.\\
\medskip
Кроме классического есть и другие определения вероятности. Их суть станет ясна немного позже:

\begin{definition}[Геометрическое определение вероятности]
Вероятность попадания случайной точки в заданную область $A$ равна отношению меры (длины, площади или объёма) этой области к мере всей возможной области исходов $\Omega$:
\[
\mathbb{P}(A) = \frac{\operatorname{mes}(A)}{\operatorname{mes}(\Omega)}.
\]
\end{definition}

\begin{definition}[Статистическое определение вероятности]
Вероятность события $A$ -- это предел, к которому стремится относительная частота наступления этого события при неограниченном увеличении числа одинаковых независимых испытаний:
\[
\mathbb{P}(A) = \lim_{n \to \infty} \frac{m}{n},
\]
где $m$ -- число появлений события, $n$ -- общее число испытаний.
\end{definition}

\begin{definition}[Аксиоматическое определение вероятности (по А.~Н.~Колмогорову)]
Вероятность -- это числовая функция (мера) $\mathbb{P}$, определённая на $\sigma$-алгебре событий пространства элементарных исходов $\Omega$, удовлетворяющая трём аксиомам:
\begin{enumerate}
    \item Неотрицательность: $\mathbb{P}(A) \ge 0$ для любого события $A$.
    \item Нормированность: $\mathbb{P}(\Omega) = 1$ (вероятность достоверного события равна 1).
    \item Счётная аддитивность: если события $A_1, A_2, \dots$ попарно несовместны, то
    \[
    \mathbb{P}\!\left(\bigcup_{i=1}^{\infty} A_i\right) = \sum_{i=1}^{\infty} \mathbb{P}(A_i).
    \]
\end{enumerate}
\end{definition}

Аксиоматическое определение является наиболее общим определением вероятности.

\end{document}